\documentclass[a4paper, 12pt]{article}
\usepackage{graphicx}
\usepackage[T2A]{fontenc} 
\usepackage[utf8]{inputenc}
\usepackage[english, russian]{babel}
\usepackage{amsmath}
\usepackage{amsfonts} 
\usepackage{amssymb}
\usepackage{xcolor}
\usepackage{array}
\usepackage[left=2cm,right=2cm,top=1cm,bottom=0.5cm,bindingoffset=0cm]{geometry}
\usepackage{float}
\usepackage{wrapfig}
\usepackage{tabularx}
\usepackage{nicematrix}
\usepackage{multirow}
\usepackage{multicol}

\begin{document}


\begin{center}
	\section*{Глава 2.}
\end{center}

\section*{3 билет. Алгоритм Евклида. Следствия из алгоритма Евклида}

\textbf{Алгоритм Евклида: общая запись.}\\

Пусть есть два числа $a,b \in \mathbb{N}, a > b$. Будем делить с остатком и при этом использовать остаток от деления первого числа для второго. Все числа здесь принадлежат $\mathbb{N}$. 
\\$r_{i}$ - это остаток, а $q_{j}$ - это целый множитель.

\begin{center}
	$a = b \cdot q_1 + r_1$\\
	$b = r_1 \cdot q_2 + r_2$\\
	$r_1 = r_2 \cdot q_3 + r_3$\\
	$\vdots$\\
	$r_{n-3} = r_{n-2} \cdot q_{n-1} + r_{n-1}$\\
	$r_{n-2} = r_{n-1} \cdot q_n + 0$\\
\end{center}

При этом целые части $q_i$ нам вообще не важны: мы используем только остатки. Основная суть этого алгоритма это использование предыдущего остатка $r_{i-1}$ от деления в следующем шаге как множитель, вместе с новым каким-то числом для получения нового, важного нам, остатка, пока не получим 0.\\

\textbf{Теорема 2.}	

Всегда последний ненулевой остаток ($r_{n-1}$) будет НОД двух чисел $a, b$.\\

Для доказательства заметим, что последовательность множителей $\{b, r_1, r_2, r_3, \cdots, r_{n-1}\}$ всегда убывает (строго). А значит мы не можем делать эти шаги бесконечно, а также последним шагом будет ноль.\\

Пусть $d$ будет общим делителем $a,b$, тогда $a \vdots d$ и $b \vdots d$. Так как $a \vdots d$, то и $(b \cdot q_1 + r_1) \vdots d$ по первому равенству. А значит левую часть можно представить как $d \cdot(\text{какое-то целое число})$. Тогда можно заметить, что остаток $r_1$ тоже должен делится на $d$, чтобы можно было вынести $d$ за скобки. Так как мы знаем, что $b$, $r_1$ делятся на $d$, то и $r_2$ будет делится на $d$, если мы проделаем все то же самое. Мы можем так продолжать до $r_{n-1}$. В общем, $OD(a,b) = OD(b,r_1) = \cdots = OD(r_{n-2}, r_{n-1}) = OD(r_{n-1}, 0) = r_{n-1}$, так как для любого $a, OD(a, 0) = a$ (можно вспомнить свойство два НОД'а). Так как у них совпадают общие делители не сложно понять, что больший из них тоже будет совпадать. Ура, получилось.\\

\textbf{Теорема 3}\\
Пусть $a, b, m \in \mathbb{N}$. Тогда $(a \cdot m, b \cdot m) = m \cdot (a, b)$.\\

При этом $m$ может принадлежать $\mathbb{Q}$, если числитель и знаменатель $\in OD(a,b)$.\\

Пример: $(14,4) = 2 \cdot (7,2) = 2 \cdot 1 = 2.$ или $(7,2) = (\frac{14}{2},\frac{4}{2}) = \frac{(14,4)}{2} = \frac{2}{2} = 1$\\

\textbf{Доказательство}

\begin{itemize}

\item Если $m \in \mathbb{N}$: 

Рассмотрим алгоритм Евклида: $a \cdot m = m \cdot (b \cdot q_1 + r_1) = b \cdot m \cdot q_1 + r_1 \cdot m$. Так как остаток $r_1$ умножается на $m$, этот множитель сохранится вплоть до $r_{n-1}$, а значит $r_{n-1} \cdot m = (a,b) \cdot m$.

\item Если $m \in \mathbb{Q}$ и $m = \frac{1}{n}, n \in \mathbb{N}$.

Так как $n \in OD(a,b)$ то число будет все равно в $\mathbb{Z}$ и мы повторяем шаги из первого случая.

\end{itemize}

\section*{4 билет. Линейное представление НОД}

\textbf{Теорема 4}\\

Пусть $a, b \in \mathbb{Z}$. Тогда существуют такие $x, y \in \mathbb{Z}$, что $(a, b) = ax + by$.
Это называется линейным представлением НОДа.

\textbf{Доказательство}.
\begin{itemize}
	
\item Сначала приведем числа к виду удобному для алгоритма Евклида. Так как делители у чисел $a$ и $-a$ одни и те же, $(a, b) = (a, -b)$. Поэтому, можно считать, что $a, b \in \mathbb{N}$.

\item НУО $a \eqslantgtr b$. Воспользуемся алгоритмом Евклида и
соответствующими обозначениями, дополним их: пусть $r_0 = b$
и $r_{-1} = a$.
\item Докажем по задне приводной индукции, для Л.П. будем брать рядом стоящие остатки.
База $k = n$: $(a, b) = 1 · r_{n-1} + 0$.\\
Переход $k \rightarrow k - 1$. Из алгоритма Евклида мы знаем, что\\
$r_{k-2} = r_{k-1} \cdot q_k + r_{k} \Rightarrow r_k = r_{k-2} - r_{k-1} \cdot q_k.$ 
\item Подставим:\\
$(a, b) = x_k \cdot r_k + y_k r_{k-1} = x_k \cdot (r_{k-2} - r_{k-1} \cdot q_k) + y_k \cdot r_{k-1} =
(-x_k \cdot q_k + y_k)r_{k-1} + x_k \cdot r_{k-2}.$

\item То есть мы перешли к предыдущему номеру остатка, карабкаясь вверх. Значит существует Л.П. для каждой пары рядом стоящих остатков, в том числе $a$ и $b$. (мы их обозначили как остатки номеров $-1$ и $0$). 
\end{itemize}

\section*{5 билет. НОД нескольких чисел через НОД двух чисел. Линейное представление НОД нескольких чисел.}

Чтобы взять НОД нескольких чисел $a_1, a_2, \hdots, a_n \in \mathbb{N}$ Нужно брать НОД чисел попарно. $(a_1, a_2, a_3) = ((a_1, a_2), a_3)$\\

Тогда мы ищем НОД только двух чисел, а по алгоритму Евклида такое действие определено.\\

Примеры:\\ 
1) $(2,4,6,8) = ((2,4), 6, 8) = (2, 6, 8) = ((2,6), 8) = (2, 8) = 2.$
\\
2) $(1,3,7) = ((1,3), 7) = (1, 7) = 1$\\


Пусть есть $a_1, a_2, \hdots a_n$, где $n > 2$. Тогда разобъем их попарно и положим в $m_i$: $m_2 = (a_1, a_2), m_3 = (a_2, a_3), \cdots, m_{n} = (a_{n-1}, a_n).$.\\

\textbf{Теорема 5}

После разбиения попарно выше, $m_n = (a_1, a_2, \hdots a_n).$ и $OD(m_n) = OD(a_1, a_2, \hdots a_n)$.\\

\textbf{Доказательство}

\begin{itemize}

\item Докажем индукцией по количеству элементов $k$.

База $k = 2$ доказана с помощью Алгоритма Евклида.

Переход $k \rightarrow k + 1$:

$OD(a_1, a_2, \hdots a_k, a_{k+1}) = OD(OD(a_1, a_2, \hdots a_k), a_{k+1}) = OD(OD(m_k), a_{k+1}) = OD(m_k, a_{k+1}) = OD(m_{k+1})$

\item Переход сработал, так что верно, что $OD(m_n) = OD(a_1, a_2, \hdots, a_n) \Rightarrow m_n = (a_1, a_2, \hdots, a_n).$
\end{itemize}

\textbf{Следствие}

Для $a_1, a_2, \hdots, a_n \in \mathbb{Z}$ существует Л.П. НОД, то есть, такие $x_1, x_2 \hdots, x_n \in \mathbb{Z}$, что $(a_1, a_2, \hdots, a_n) = x_1 \cdot a_1 + x_2 \cdot a_2 + \cdots + x_n \cdot a_n$.

\textbf{Доказательство}

\begin{itemize}
	
\item Докажем индукцией по количеству элементов $k$.

База $k = 2$ Доказана в Теореме 4. (Л.П. двух чисел).

Переход $k \rightarrow k + 1$:

Воспользуемся Теоремой 5.

$(a_1, a_2, \hdots, a_k, a_{k+1}) = (m_k, a_{k+1}).$

Представим первую скобку в виде Л.П. для $y$ и $x_{k+1}$:

$(m_k, a_{k+1}) = y \cdot m_k + x_{k+1} \cdot a_{k+1}$\\

Раскроем $m_k$ по ее определению из Теоремы 5 и применим индукционное предположение, где множителями будут $x_1, x_2, \hdots, x_k.$

$(a_1, a_2, \hdots, a_{k+1}) =  y \cdot m_k + x_{k+1} \cdot a_{k+1} = y \cdot (x_1 \cdot a_1 + x_2 \cdot a_2 + \cdots + x_k \cdot a_k) + x_{k+1} \cdot a_{k+1}$

Раскроем скобки

$(a_1, a_2, \hdots, a_{k+1}) =  y \cdot m_k + x_{k+1} \cdot a_{k+1} = (y \cdot x_1) \cdot a_1 + (y \cdot x_2) \cdot a_2 + \cdots + (y \cdot x_k) \cdot a_k + x_{k+1} \cdot a_{k+1}$

\item Сейчас в скобках нужные нам коэффициенты для Л.П. и, так как сработал индукционный переход, существует Л.П. для НОДа сразу нескольких чисел.

\end{itemize}

\end{document}
